\documentclass[11pt]{article}
\usepackage[a4paper, margin=2.54cm]{geometry}

% español
\usepackage[spanish]{babel}

% imágenes
\usepackage{graphicx}
%\graphicspath{{img}}

% fuentes de conjuntos numéricos
\usepackage{amsfonts}

% símbolos
\usepackage{amsmath, amssymb}

% gráficos
%\usepackage{tikz}

% plots
%\usepackage{pgfplots}
%\pgfplotsset{width=10cm, compat=1.9}

% shapes
%\usetikzlibrary{shapes.geometric}

% espacio entre párrafos
\usepackage{parskip}

% cancelar términos en ecuaciones
%\usepackage{cancel}

% links
%\usepackage[colorlinks=true,
%    urlcolor=blue]{hyperref}

% incluir pdfs
%\usepackage{pdfpages}

% tipografía cheta
%\usepackage{tgadventor}
%\renewcommand*\familydefault{\sfdefault}
%\usepackage[T1]{fontenc}

% interlineado: 1.0 simple, 1.3 (1.5), 1.6 (doble)
%\linespread{1.3}

% bibliografía
%\usepackage[backend=biber, style=ieee, sorting=nyt]{biblatex}
%\bibliography{bib}
%\usepackage{csquotes}

\begin{document}

\begin{titlepage}
    \begin{center}
        \vspace*{.5cm}
        \includegraphics[scale=.5]{../../../img/udemm-logo.png}\\
        \vspace{.2cm}
        \Large
        \textbf{Facultad de Ingeniería}\\
        \textbf{Ingeniería en Sistemas}\\
        \vspace{2cm}

        \Huge
        Diseño de Sistemas II \\
        Trabajo Práctico \\ 
        Arquitectura Cliente/Servidor\\
        \vfill

        \raggedright
        \Large
        Docente: Mg. Daniel Omar García \\
        Alumno: Daniel Ise \\
        Legajo: 28547 \\
        Fecha: Octubre, 2026 
    \end{center}
\end{titlepage}

\newpage

\tableofcontents
\newpage

\section{Introducción}

La arquitectura Cliente-Servidor es un modelo de diseño de \emph{Software},
basado en la división del sistema en dos componentes principales:
el \textbf{Cliente} y el \textbf{Servidor}.

Su funcionamiento sigue el \emph{ciclo de petición y respuesta}:
\begin{itemize}
    \item El Cliente envía una \emph{petición} al Servidor,
    solicitando un recurso.

    \item El Servidor \emph{procesa} la solicitud y elabora una
    \emph{respuesta}.
    
    \item El Servidor \emph{envía} esa respuesta al Cliente.

    \item El Cliente recibe la respuesta y la procesa.
\end{itemize}

Toda la comunicación sigue un \textbf{protocolo},
que especifica las \emph{reglas} que tanto el Servidor como el Cliente usan para
comunicarse.

Por ejemplo, el navegador web \emph{Chromium} envía una petición a un servidor
web, servido con \emph{Apache}.
El servidor recibe la solicitud, la procesa y elabora una respuesta, que puede
ser el renderizado de una página web escrita en \emph{PHP}.
La respuesta se envía al cliente, que la renderiza para el usuario.
Toda la comunicación se hace siguiendo las reglas especificadas por los
protocolos TCP/IP y HTTP.

La arquitectura Cliente-Servidor cuenta entre sus ventajas:
\begin{itemize}
    \item La \emph{centralización} de los recursos y los datos en un Servidor.
    
    \item La posibilidad de \emph{administrar la información}desde el Servidor.
    
    \item La posibilidad de \emph{escalar} el servicio añadiendo más clientes.
    
    \item El \emph{control} de la seguridad desde el Servidor.
\end{itemize}

Aunque cuenta con algunas desventajas:
\begin{itemize}
    \item \emph{Dependencia del servidor}:
    ante una falla, los clientes pierden acceso al servicio.
    
    \item Puede requerir de una \emph{inversión importante} en Hardware y
    Seguridad 
    
    \item Puede enfrentar problemas de \emph{saturación} frente a picos de 
    clientes en simultáneo.
\end{itemize}

Para el trabajo práctico se eligieron cuatro casos exitosos de implementación de 
esta arquitectura:
\begin{itemize}
    \item Netflix, servicio de \emph{streaming} de video.
    
    \item Spotify, servicio de \emph{streaming} de música y \emph{Podcasts}.

    \item Whatsapp, aplicación de mensajería instantánea.
    
    \item Gmail, aplicación pionera en la provisión de \emph{email} en el
    navegador.
\end{itemize}

Se eligieron estas aplicaciones por ser casos exitosos en los desafíos más 
importantes que enfrentan este tipo de arquitecturas.
Todas deben recibir, procesar y enviar grandes volúmenes de datos.
Deben ser capaces de responder a grandes picos de demanda.
Están obligadas a cuidar la seguridad y privacidad de sus clientes,
ofreciendo servicios de autenticación y autorización.
Todo ello mientras tienen presencia en distintos dispositivos,
incluyendo computadoras de escritorio, teléfonos, portátiles y televisores.
En otras palabras, todas enfrentan grandes desafíos de rendimiento,
escalabilidad, seguridad y experiencia de usuario.

En lo que respecta a los casos de implementaciones problemáticas, se eligieron:
\begin{enumerate}
    \item SIU Guaraní
    \item El lanzamiento del videojuego SimCity
    \item Bolsas y Mercados Argentinos (BYMA)
\end{enumerate}

En los tres casos se observan problemas comunes en la arquitectura de
Cliente-Servidor: cuellos de botella, problemas de seguridad, dependencia
crítica y sobrecarga del servidor.

\section{Casos exitosos}

\subsection{Netflix}

Netflix es un proveedor de servicios de \emph{streaming} de video, enfocado en
series y películas.
Enfrenta grandes desafíos de ingeniería en lo que respecta a rendimiento:
debe recibir, procesar y enviar grandes volúmenes de datos a cientos de miles de
clientes de todo el mundo, que se conectan desde dispositivos muy diferentes en
capacidades de cómputo.
El Servidor debe de proveer autenticación, catálogo, perfiles y transmisión de
video en tiempo real, por lo que supone un caso de análisis particularmente
interesante.

\subsection{Spotify}

El segundo caso a considerar es el de Spotify. Se trata también de un servicio 
de \emph{streaming}, en este caso de audio, siendo quizá el servicio más
exitoso en su rubro. Sus servidores deben de proveer a los clientes, que
utilizan todo tipo de dispositivos, de autenticación, recomendaciones, acceso a 
catálogo y transmisión de audio. Se trata de otro caso de una empresa de gran 
éxito, con una amplia variedad de clientes (en términos de dispositivos
diferenciados), por lo que su análisis reviste de particular importancia.

\subsection{Whatsapp}

El tercer caso de éxito a considerar es el de \textbf{Whatsapp}, una aplicación
de mensajería instantánea, orientada principal aunque no exclusivamente a la 
mensajería móvil. El servidor debe de gestionar colas de mensajes,
notificaciones, autenticación y encriptación, tratando de minimizar, en los
casos en que corresponda, el consumo de batería en sus clientes. La gran
difusión de esta aplicación, así como el desafío de ingeniería que supone el
manejo del enorme flujo de datos diarios, motiva a incluirla en este trabajo 
práctico.

\subsection{Gmail}

Por último, el cuarto caso es el de \textbf{Gmail}, aplicación pionera en la 
provisión de \emph{email} en navegador, que utiliza tecnología AJAX del lado del 
cliente, para permitir el uso del correo electrónico de forma eficiente sin 
requerir de la instalación de aplicaciones de escritorio, como era el estándar 
del rubro en el momento de su aparición. Su desarrollo y evolución como 
aplicación que implementa la arquitectura de Cliente-Servidor parece casi obvia,
por la innovación que representó en su momento y por ser pionera en las
aplicaciones web asíncronas, modelo para las que vendrían luego.

\section{Casos problemáticos}

\subsection{SIU Guaraní}

\subsection{SimCity}

\subsection{Bolsas y Mercados Argentinos (BYMA)}

\section{Conclusiones}

%\printbibliography[heading=bibintoc, title={Bibliografía}]


\end{document}
